\chapter{\textbf{CONSIDERAÇÕES FINAIS}}

Este capítulo consolida as principais conclusões do estudo, discute suas limitações e aponta direções para pesquisas futuras que possam vir a cobrir as lacunas deixadas pela presente pesquisa.

O presente trabalho investigou o impacto da aplicação do critério de validação de tábuas de mortalidade estabelecido pelo Art. 36, inciso I, da Portaria MTP nº 1.467/2022, sobre as PM dos RPPS. O critério em questão determina que, caso o ente não use a tábua de referência, a tábua adotada deve apresentar, na idade média da massa segurada, uma expectativa de vida igual ou superior àquela calculada pela tábua de referência. Em todos os quatro perfis etários analisados foram identificadas tábuas de mortalidade que atenderam ao critério normativo, isto é, apresentaram expectativa de vida na idade média superior à da tábua de referência, mas que, ao mesmo tempo, produziram uma PM inferior àquela gerada pela tábua de referência (IBGE 2024 Extrapolada para cada gênero).

As tábuas que apresentaram uma PM inferior variaram conforme o perfil etário e o método de custeio adotado. Com relação ao perfil etário, as simulações indicam que massas seguradas com estrutura etária mais jovem, especialmente a do sexo masculino, podem ser mais suscetíveis à ocorrência de inconsistências. Na base jovem, onde aproximadamente 60\% dos segurados concentram-se na faixa de 18 a 30 anos, os métodos INE e CUP identificaram dez e oito inconsistências para o sexo masculino, respectivamente, além de três para o sexo feminino em cada um desses métodos. O método de custeio PNI, por sua vez, apresentou o maior número total de inconsistências em todos os perfis, somando 16 ocorrências na base jovem, sendo 12 para o sexo masculino e 4 para o sexo feminino. A base equilibrada mostrou um comportamento semelhante a base jovem, embora tenha apresentando um numero menor de inconsistências.

Em contrapartida, a base idosa, com concentração majoritária de segurados acima dos 51 anos e idade média de 65,13 anos, apresentou o menor número de inconsistências nos métodos INE e CUP, registrando apenas uma tábua inconsistente para cada sexo. A base bimodal comportou-se de maneira semelhante à base idosa nesses dois métodos, em razão do elevado peso relativo dos grupos de aposentados.

A investigação do comportamento da curva de $l_x$ indicou que a assimetria entre os níveis de sobrevivência nas tábuas pode ser o motivo de tais inconsistências. Uma tábua comparativa pode apresentar níveis de sobrevivência superiores à tábua de referência ao longo das idades contributivas $[18, 64]$, o que eleva o VACF, e simultaneamente apresentar níveis de sobrevivência inferiores nas idades de aposentadoria $[65, \omega]$, reduzindo o VABF. Tal assimetria explica o aumento da expectativa de vida na tábua comparativa, além de implicar na variação do VACF superando a variação no VABF, resultando em uma PM inferior à da tábua de referência, ainda que a expectativa de vida na idade média seja igual ou superior.

Esse mecanismo foi quantificado pelo índice $I_F$ proposto na Equação \ref{eq:indice}, o qual tenta mensurar o diferencial de sobrevivência ponderado pela distribuição etária da massa, separadamente para as faixas contributivas e para as faixas em gozo de benefício. A análise dos índices para as tábuas inconsistentes das bases jovem e equilibrada evidenciou que, em sua maioria, essas tábuas apresentam $I_{Ativos} > I_{Aposentados}$, confirmando o padrão esperado. As exceções observadas, em particular no método PNI, reforçam que o índice proposto, embora útil como instrumento de diagnóstico prévio, não incorpora a complexidade total no cálculo da PM, principalmente no método de custeio PNI.

Adicionalmente, os resultados indicaram uma possível vulnerabilidade normativa específica. O mecanismo de suavização e agravamento de tábuas de mortalidade pode ser utilizado para modelar tábuas que, embora tecnicamente aprovadas pelo critério do Art. 36, produzem PM inferiores às da tábua de referência. Tábuas como a BR-EMSmt-v.2021 M Agravada em 24\% e Prudential-50 Suavizada apareceram de forma consistente em múltiplos perfis e métodos, embora suas versões sem alterações não aparecem. A esse respeito, a pesquisa sugere que é possível suavizar ou agravar uma tábua de mortalidade de tal forma que ela mantenha a expectativa de vida na idade média próxima à da tábua de referência, satisfazendo o critério normativo, mas resultando em uma PM inferior.

Em síntese, o estudo indica que o critério de validação estabelecido pelo Art. 36 da Portaria MTP nº 1.467/2022, ao utilizar a expectativa de vida na idade média como único parâmetro de comparação, pode ser insuficiente para garantir, em todas as configurações demográficas possíveis, que a tábua adotada pelo RPPS resulte em uma provisão matemática igual ou superior àquela gerada pela tábua do IBGE.

Embora os resultados obtidos sejam consistentes, algumas limitações devem ser consideradas na interpretação deste trabalho. A primeira limitação diz respeito ao uso exclusivo de dados simulados. A construção de bases populacionais sintéticas, embora seja fundamental para isolar o efeito da estrutura etária sobre o passivo atuarial em condições controladas, impede a extrapolação direta dos resultados quantitativos para RPPS específicos. As magnitudes das inconsistências identificadas, bem como o número exato de tábuas inconsistentes por perfil, são sensíveis às configurações das massas simuladas e às premissas adotadas. A aplicação da metodologia a bases reais de RPPS poderia revelar padrões distintos em termos de intensidade e frequência das inconsistências.

A segunda limitação está relacionada à simplificação das variáveis consideradas na análise. Para evidenciar com maior clareza o impacto das tábuas de mortalidade sobre a PM, o estudo restringiu-se ao benefício de aposentadoria por idade, desconsiderando outros decrementos relevantes no contexto previdenciário, como a invalidez e a reversão de pensão, por exemplo. Da mesma forma, adotou-se valor único de remuneração inicial ($R\$~5.000,00$) para todos os segurados, abstraindo a heterogeneidade salarial presente nas massas reais dos RPPS. Essas simplificações, ao reduzirem a complexidade do modelo, podem subestimar ou superestimar o efeito das inconsistências em cenários mais aderentes à realidade.

Por fim, o índice $I_F$ proposto, embora útil como instrumento de diagnóstico e interpretação das inconsistências, não incorpora explicitamente o efeito do crescimento salarial sobre o VACF no método PNI. Isso implica que o índice, em sua forma atual, possui menor poder preditivo para inconsistências nesse método específico, representando uma limitação da ferramenta diagnóstica desenvolvida neste trabalho.

Uma primeira direção natural consiste na replicação da metodologia aqui desenvolvida com dados reais de RPPS. A aplicação do índice $I_F$ a bases reais permitiria avaliar a prevalência das inconsistências identificadas nas simulações no contexto factual dos regimes brasileiros, além de verificar se determinados perfis demográficos observados na prática são mais expostos ao problema normativo descrito.

Uma segunda linha de pesquisa poderia estender o modelo atuarial para incorporar múltiplos decrementos, contemplando invalidez, desligamento e pensão por morte. A inclusão dessas variáveis tornaria o modelo mais próximo da realidade dos RPPS e permitiria avaliar se o fenômeno das inconsistências se mantém ou se atenua quando a interação entre diferentes riscos biométricos é considerada.

Em terceiro lugar, seria relevante investigar a sensibilidade das inconsistências em relação às premissas econômicas, particularmente à taxa de juros e à taxa de crescimento salarial. Análises de cenários com diferentes trajetórias dessas variáveis contribuiriam para mapear as condições sob as quais as inconsistências são mais ou menos graves.

Outra perspectiva é o aprimoramento do índice $I_F$, de modo a incorporar a complexidade do cálculo da PM em cada método de custeio. Um índice estendido que contemplasse, simultaneamente, os diferenciais de sobrevivência, de crescimento salarial e de outras premissas atuariais, apresentaria maior poder preditivo, tornando-se uma ferramenta robusta para o diagnóstico prévio de inconsistências.

Do ponto de vista regulatório, estudos futuros poderiam propor e testar critérios alternativos ou complementares ao estabelecido no Art. 36 da Portaria MTP nº 1.467/2022. Entre as possibilidades, sugerem-se critérios baseados em métricas distribucionais das tábuas de mortalidade como percentis ponderados pela estrutura etária da massa. Tal critério poderia ser mais robustos a assimetrias demográficas do que a comparação pontual de expectativas de vida.

Por fim, pesquisas que examinem o comportamento de tábuas geracionais no contexto do Art. 36 representam uma extensão importante. As tábuas geracionais incorporam a perspectiva de melhoria contínua da mortalidade ao longo do tempo e são reconhecidas na literatura como mais adequadas para a mensuração do risco de longevidade em planos de longo prazo. A investigação de como o critério normativo se comporta diante dessas tábuas, especialmente em massas jovens, poderia revelar novas dimensões do problema identificado neste estudo.